problem:  
Let \( M^{2n} \) be a simply connected, compact Riemannian manifold with nonnegative sectional curvature. Suppose \( M \) admits an effective isometric torus action of rank  
\[
\mathrm{rk}(T^k) \ge \frac{n}{2},
\]
and that the fixed point set of this action is nonempty.  

**Question:**  
Must \( M \) be *rationally elliptic*, i.e., does \( M \) have only finitely many nontrivial rational homotopy groups?  
Equivalently, can the combination of nonnegative sectional curvature and large torus symmetry (of rank at least \( n/2 \)) force finiteness of rational homotopy, as predicted by the Bott conjecture?  
If so, what is the geometric or topological mechanism—through curvature, orbit-type stratification, or fixed-point data—that enforces rational ellipticity under these symmetry constraints?  

Why is it a "good" problem:  
This problem formulates a precise intermediate step toward the Bott conjecture: instead of tackling the general case, it isolates the scenario where a compact, simply connected, nonnegatively curved manifold has symmetry rank at least half its dimension. This symmetry condition is motivated by known rigidity phenomena—high symmetry tends to severely restrict manifold topology and curvature behavior (as shown in the Grove–Searle symmetry-rank theorems). Testing whether this degree of symmetry guarantees rational ellipticity creates a natural quantitative frontier between known results (where rational ellipticity is proved for maximal or almost-maximal symmetry) and the full conjecture (where no symmetry is assumed). The question thus sits at a conceptual crossroads linking Riemannian geometry, transformation group theory, and rational homotopy theory, offering a sharply posed and potentially generative step toward deeper structural understanding. It is mathematically precise, logically sound, and opens multiple avenues of attack—geometric, homotopical, and equivariant—making it an exemplary “narrow entrance, profound depth” research problem.